%!TEX TS-program = xelatex
\documentclass[fontset=none]{ctexart}
\usepackage[a5paper,margin=20mm]{geometry}
\pagestyle{empty}
\setCJKmainfont[BoldFont=FandolHei-Regular.otf]{FandolSong-Regular.otf}

% Maple Mono NF CN can be installed system-wide, or its Regular TTF can be
% placed next to this MWE after downloading the upstream release archive.
\IfFontExistsTF{Maple Mono NF CN}
  {
    \setCJKmonofont{Maple Mono NF CN}
    \setmonofont{Maple Mono NF CN}
  }
  {
    \setCJKmonofont[Path=./]{MapleMono-NF-CN-Regular.ttf}
    \setmonofont[Path=./]{MapleMono-NF-CN-Regular.ttf}
  }

% Select the monospaced family first, then apply xeCJK's existing code-layout
% model inside the command's local group.
\DeclareTextFontCommand{\InlineCode}{\ttfamily\xeCJKVerbAddon}

\begin{document}
\section*{Issue \#808：行内代码的中西文间距}

\texttt{\string\texttt} 只是字体切换，仍按普通正文规则在中西文之间
加入间距；语义化的行内代码命令还需要启用 \texttt{\string\xeCJKVerbAddon}。

\bigskip
{\Large
\noindent 默认：开始\texttt{TEST测试TEST}结束\par
\medskip
\noindent 建议：开始\InlineCode{TEST测试TEST}结束\par
\medskip
\noindent 混排：开始\InlineCode{TE测STTE试ST}结束\par
}

\bigskip
建议版本在代码内部保持半角/全角网格，同时保留代码片段与正文之间的
正常边界间距。该模式会禁止片段内部自动断行，因此不应无条件应用到
所有 \texttt{\string\ttfamily} 文字。
\end{document}
